% Trascrizione del PDF quiz Paola Suria-2-061-070(2).pdf
% Pagine 061, 067-069: spiegazioni, non quiz.
% Le soluzioni 13a simulazione provengono dalla tabella a pagina 066.
% La pagina 070 non contiene una tabella delle soluzioni.

% Pag. originale 62/521
\iffalse Pag. originale 62/521 \fi
\begin{quiz}
Quale delle seguenti funzioni è pari?

\begin{enumerate}
\item $f(x)=e^x\sin 2x$
\item $f(x)=x\cos x^3$
\item $f(x)=e^{|x|}\sin(x^3+x)$
\item $f(x)=(x^3+1)\cos 2x$
\item nessuna risposta è esatta
\end{enumerate}
\end{quiz}

% SOLUZIONE: E
% SEZIONE_AM1: 1.4
% SPIEGAZIONE: Le prime quattro funzioni non sono pari; la risposta ufficiale è E.

% Pag. originale 62/521
\iffalse Pag. originale 62/521 \fi
\begin{quiz}
I valori $\alpha,\beta$ per i quali $f(x)=\begin{cases}\alpha x^2+\beta&x\le1\\2\beta x^2-\alpha x+1&x>1\end{cases}$ è derivabile in $(-\infty,+\infty)$

\begin{enumerate}
\item $\alpha=-3,\ \beta=-4$
\item $\alpha=-4,\ \beta=-3$
\item $\alpha=\frac35,\ \beta=\frac45$
\item $\alpha=\frac45,\ \beta=\frac35$
\item $\alpha=\frac12,\ \beta=0$
\end{enumerate}
\end{quiz}

% SOLUZIONE: D
% SEZIONE_AM1: 4.1
% SPIEGAZIONE: Imponendo continuità e uguaglianza delle derivate in $x=1$ si ottiene $\alpha=4/5$, $\beta=3/5$.

% Pag. originale 62/521
\iffalse Pag. originale 62/521 \fi
\begin{quiz}
$\displaystyle\lim_{x\to0}\frac{x-\log(1+x)}{x\sin 3x}=$

\begin{enumerate}
\item $\frac16$
\item $4$
\item $-4$
\item $\frac14$
\item $\frac13$
\end{enumerate}
\end{quiz}

% SOLUZIONE: A
% SEZIONE_AM1: 3.1
% SPIEGAZIONE: Il numeratore è equivalente a $x^2/2$ e il denominatore a $3x^2$; il limite è $1/6$.

% Pag. originale 62/521
\iffalse Pag. originale 62/521 \fi
\begin{quiz}
L’insieme dei numeri $\{z\in\mathbb C:\operatorname{Re}(z\overline z)=1\}$ sul piano di Argand Gauss è rappresentato da:

\begin{enumerate}
\item due rette
\item un’iperbole
\item un cerchio
\item una circonferenza
\item vuoto
\end{enumerate}
\end{quiz}

% SOLUZIONE: D
% SEZIONE_AM1: 1.5
% SPIEGAZIONE: Poiché $z\overline{z}=|z|^2=1$, il luogo è la circonferenza unitaria.

% Pag. originale 62/521
\iffalse Pag. originale 62/521 \fi
\begin{quiz}
Se $g:\mathbb R\to\mathbb R$ una funzione derivabile una volta, allora integrando per parti $\displaystyle\int_0^2 xg'(x)\,dx$ si ottiene:

\begin{enumerate}
\item $g(2)-g(0)$
\item $g(2)+g(0)$
\item $\int_0^2 g(x)\,dx$
\item $\left[\frac{x^2}{2}g\prime(x)\right]_0^2-\int_0^2\frac{x^2}{2}g\prime\prime(x)\,dx$
\item $2g(2)-\int_0^2 g(x)\,dx$
\end{enumerate}
\end{quiz}

% SOLUZIONE: E
% SEZIONE_AM1: 5.4
% SPIEGAZIONE: Integrazione per parti con $u=x$ e $dv=g\prime(x)\,dx$: si ottiene $2g(2)$ meno l’integrale di $g$ sugli estremi indicati nel quesito.

% Pag. originale 63/521
\iffalse Pag. originale 63/521 \fi
\begin{quiz}
Se $f:(0,+\infty)\to\mathbb R$ definita da $f(x)=\log x+x$. Allora per $x\to+\infty$

\begin{enumerate}
\item l’asintoto obliquo è la semiretta $y=x+1$
\item l’asintoto obliquo non esiste
\item $f(x)=x+o(1)$
\item l’asintoto obliquo è la semiretta $y=x$
\item l’asintoto obliquo è la semiretta $y=x+e$
\end{enumerate}
\end{quiz}

% SOLUZIONE: B
% SEZIONE_AM1: 4.4
% SPIEGAZIONE: $f(x)-x=\log x$ diverge, quindi non esiste intercetta finita per un asintoto obliquo.

% Pag. originale 63/521
\iffalse Pag. originale 63/521 \fi
\begin{quiz}
Sia $g:\mathbb R\to\mathbb R$: $\displaystyle\int_0^2 g(3t+2)\,dt=$

\begin{enumerate}
\item $g(8)-g(2)$
\item $2g(5)$
\item $\frac13\int_2^8 g(t)\,dt$
\item $3\int_2^8 g(t)\,dt$
\item $\frac13\int_0^2 g(t)\,dt$
\end{enumerate}
\end{quiz}

% SOLUZIONE: C
% SEZIONE_AM1: 5.4
% SPIEGAZIONE: Sostituendo $u=3t+2$ si ha $dt=du/3$, con estremi $2$ e $8$.

% Pag. originale 63/521
\iffalse Pag. originale 63/521 \fi
\begin{quiz}
Sia $f:\mathbb R\to\mathbb R$ continua. Allora $\displaystyle\int_{-1}^{1} f(x^2)\,dx=$

\begin{enumerate}
\item $2\int_0^1 f(x^2)\,dx$
\item $\int_0^1\frac{f(t)}{2t}\,dt$
\item $\int_{-1}^1 f(t)2t\,dt$
\item $2\int_{-1}^0 f(x^2)\,dx$
\item $0$
\end{enumerate}
\end{quiz}

% SOLUZIONE: E
% SEZIONE_AM1: 5.2
% SPIEGAZIONE: DISCREPANZA: risposta ufficiale E, ma integrando una funzione pari le opzioni A e D sono vere; E è falsa in generale (esempio $f=1$).

% Pag. originale 63/521
\iffalse Pag. originale 63/521 \fi
\begin{quiz}
L’integrale improprio $\displaystyle\int_0^{+\infty}\frac{dt}{t^\alpha+t^{2\alpha}}$ converge se

\begin{enumerate}
\item $\frac12<\alpha<1$
\item $1<\alpha$
\item nessun valore di $\alpha$
\item $\alpha<\frac12$
\item $\forall\alpha$
\end{enumerate}
\end{quiz}

% SOLUZIONE: A
% SEZIONE_AM1: 5.5
% SPIEGAZIONE: Per $\alpha>0$, vicino a zero serve $\alpha<1$ e all’infinito serve $2\alpha>1$.

% Pag. originale 63/521
\iffalse Pag. originale 63/521 \fi
\begin{quiz}
Il polinomio di Mac Laurin di grado 2 della funzione $f(x)=e^{x^2+2x}$ è:

\begin{enumerate}
\item $1+2x+4x^2$
\item $1+2x+4x^2$
\item $2x+4x^2$
\item $1+2x+3x^2$
\item $1+2x+5x^2$
\end{enumerate}
\end{quiz}

% SOLUZIONE: D
% SEZIONE_AM1: 4.1
% SPIEGAZIONE: Espandendo $e^u=1+u+\frac{u^2}{2}+o(u^2)$, con $u=2x+x^2$, si ottiene $1+2x+3x^2+o(x^2)$. A e B sono duplicate nel PDF.

% Pag. originale 64/521
\iffalse Pag. originale 64/521 \fi
\begin{quiz}
Sia $f$ una funzione continua e positiva $\forall x\in\mathbb R$. Posto $F(x)=\displaystyle\int_0^x f(t)\,dt$ segue necessariamente che

\begin{enumerate}
\item $\lim_{x\to+\infty}F(x)$ è finito
\item $F(-1)<0$
\item $F$ ha un minimo in $\mathbb R$
\item $\lim_{x\to+\infty}F(x)=+\infty$
\item $\lim_{x\to+\infty}f(x)=+\infty$
\end{enumerate}
\end{quiz}

% SOLUZIONE: B
% SEZIONE_AM1: 5.2
% SPIEGAZIONE: $F$ è strettamente crescente e $F(0)=0$, perciò $F(-1)<0$.

% Pag. originale 64/521
\iffalse Pag. originale 64/521 \fi
\begin{quiz}
$\displaystyle\int_0^1 f(2x)\,dx=$

\begin{enumerate}
\item $2\int_0^2 f(x)\,dx$
\item $2\int_0^{1/2} f(x)\,dx$
\item $2\int_0^1 f(x)\,dx$
\item $\frac12\int_0^1 f(t)\,dt$
\item $\frac12\int_0^2 f(t)\,dt$
\end{enumerate}
\end{quiz}

% SOLUZIONE: E
% SEZIONE_AM1: 5.4
% SPIEGAZIONE: Con $t=2x$, $dx=dt/2$ e gli estremi diventano $0$ e $2$.

% Pag. originale 64/521
\iffalse Pag. originale 64/521 \fi
\begin{quiz}
Sia $a>0$. Allora $\displaystyle\int_a^{+\infty}x^{-\frac43}\,dx=$

\begin{enumerate}
\item $3\sqrt[3]{a}$
\item $-3\sqrt[3]{a}$
\item $+\infty$
\item $\frac{3}{\sqrt[3]{a}}$
\item $-\frac{3}{\sqrt[3]{a}}$
\end{enumerate}
\end{quiz}

% SOLUZIONE: D
% SEZIONE_AM1: 5.5
% SPIEGAZIONE: Una primitiva è $-3x^{-1/3}$; il limite superiore vale zero e resta $3/a^{1/3}$.

% Pag. originale 64/521
\iffalse Pag. originale 64/521 \fi
\begin{quiz}
Lo sviluppo di Mac Laurin di grado 3 della funzione $\log(1+x+x^2)$ è:

\begin{enumerate}
\item $x+\frac{x^2}{2}-\frac23x^3+o(x^3)$
\item $x-\frac{x^2}{2}+\frac{x^3}{3}+o(x^3)$
\item $x+x^2+x^3+o(x^3)$
\item $x+\frac{x^2}{2}+x^3+o(x^3)$
\item $x+\frac{x^2}{2}+\frac23x^3+o(x^3)$
\end{enumerate}
\end{quiz}

% SOLUZIONE: A
% SEZIONE_AM1: 4.1
% SPIEGAZIONE: Si usa $\log(1+u)=u-\frac{u^2}{2}+\frac{u^3}{3}+o(u^3)$, con $u=x+x^2$.

% Pag. originale 64/521
\iffalse Pag. originale 64/521 \fi
\begin{quiz}
L’integrale $I=\displaystyle\int_0^1\frac{\sin t}{t^\alpha}\,dt=+\infty$ se:

\begin{enumerate}
\item $\alpha\ge-1$
\item $\alpha<2$
\item $\alpha\ge2$
\item $\alpha<-1$
\item $\alpha\le1$
\end{enumerate}
\end{quiz}

% SOLUZIONE: C
% SEZIONE_AM1: 5.5
% SPIEGAZIONE: Vicino a zero $\sin t/t^{\alpha}$ è equivalente a $t^{1-\alpha}$; diverge se $\alpha\geq 2$.

% Pag. originale 65/521
\iffalse Pag. originale 65/521 \fi
\begin{quiz}
$\displaystyle\int_{-\infty}^{+\infty}\frac{dt}{1+t^2}=$

\begin{enumerate}
\item $\pi$
\item $\frac\pi2$
\item $+\infty$
\item $0$
\item $2\pi$
\end{enumerate}
\end{quiz}

% SOLUZIONE: A
% SEZIONE_AM1: 5.5
% SPIEGAZIONE: La primitiva è $\arctan t$ e il salto fra gli estremi infiniti vale $\pi$.

% Pag. originale 65/521
\iffalse Pag. originale 65/521 \fi
\begin{quiz}
Sia data la funzione $f:\mathbb R\to\mathbb R$, derivabile almeno una volta, allora: $\displaystyle\int_1^2\frac{f(x)}{x}\,dx=$

\begin{enumerate}
\item $\log2\cdot f(2)-\int_1^2\log x\cdot f\prime(x)\,dx$
\item $\log2\cdot f(2)-\int_1^2\log|x|\cdot f\prime(x)\,dx$
\item $\log2\cdot f\prime(2)-\int_1^2\log x\cdot f(x)\,dx$
\item $\frac{f(2)}2-\int_1^2\frac{f\prime(x)}x\,dx$
\item $\frac{f(2)}2-f(1)-\int_1^2\frac{f\prime(x)}x\,dx$
\end{enumerate}
\end{quiz}

% SOLUZIONE: A
% SEZIONE_AM1: 5.4
% SPIEGAZIONE: Per parti con $dv=dx/x$. DISCREPANZA: A e B sono equivalenti su $[1,2]$; la tabella ufficiale indica A.

% Pag. originale 65/521
\iffalse Pag. originale 65/521 \fi
\begin{quiz}
Sia $f:[a,b]\to\mathbb R$, continua e sia $m=\displaystyle\frac1{b-a}\int_a^b f(x)\,dx$ il valor medio integrale della funzione $f$ in $[a,b]$. Quale delle seguenti affermazioni è corretta?

\begin{enumerate}
\item Se $2<f(x)<4\ \forall x\in[a,b]$, allora $2<m<4$
\item $m=\frac{f(b)-f(a)}{b-a}$
\item Se $m<0$ allora $f(x)<0\ \forall x\in[a,b]$
\item Se $2<x<4$ allora $2<m<4$
\item $m=\frac{f(b)+f(a)}{b-a}$
\end{enumerate}
\end{quiz}

% SOLUZIONE: A
% SEZIONE_AM1: 5.2
% SPIEGAZIONE: Integrando le disuguaglianze strette $2<f<4$ e dividendo per $b-a>0$ si ottiene $2<m<4$.

% Pag. originale 65/521
\iffalse Pag. originale 65/521 \fi
\begin{quiz}
$\displaystyle\int_1^2\frac{f(\log x)}x\,dx=$

\begin{enumerate}
\item $\int_1^2 f(t)\,dt$
\item $\int_1^2\frac{f(\log x)}x\,dx$
\item $\int_1^2\frac{f(t)}{e^t}\,dt$
\item $\int_0^{\log2}f(t)\,dt$
\item nessuna delle altre risposta è corretta
\end{enumerate}
\end{quiz}

% SOLUZIONE: D
% SEZIONE_AM1: 5.4
% SPIEGAZIONE: Con $t=\log x$ e $dt=dx/x$ gli estremi diventano $0$ e $\log 2$. DISCREPANZA: anche B ripete identicamente il membro sinistro.

% Pag. originale 66/521
\iffalse Pag. originale 66/521 \fi
\begin{quiz}
$\displaystyle\int_1^{+\infty}x^\alpha\arctan x\,dx$ converge se:

\begin{enumerate}
\item $\alpha\in(-1,+\infty)$
\item $\alpha\in(-\infty,-1)$
\item $\alpha\in(-\infty,+\infty)$
\item $\alpha\in(-\infty,0)$
\item per nessuna $\alpha$
\end{enumerate}
\end{quiz}

% SOLUZIONE: B
% SEZIONE_AM1: 5.5
% SPIEGAZIONE: Poiché $\arctan x$ tende a $\pi/2$, il comportamento è quello di $x^{\alpha}$; convergenza per $\alpha<-1$.

% Pag. originale 70/521
\iffalse Pag. originale 70/521 \fi
\begin{quiz}
Per quali valori reali di $x$ le due funzioni $f(x)=\sqrt{\frac{x+1}{x-1}}$ e $g(x)=\frac{\sqrt{x+1}}{\sqrt{x-1}}$ coincidono?

\begin{enumerate}
\item $\forall x\in\mathbb R$
\item $\forall x\in(1,+\infty)$
\item $\forall x\in[1,+\infty)$
\item $\forall x\in(-\infty,-1)$
\item $\forall x\in[-1,1)$
\end{enumerate}
\end{quiz}

% SOLUZIONE: B
% SEZIONE_AM1: 1.4
% SPIEGAZIONE: Il dominio comune è $x>1$; su tale dominio vale la proprietà della radice del quoziente.

% Pag. originale 70/521
\iffalse Pag. originale 70/521 \fi
\begin{quiz}
Lo sviluppo del terz’ordine della funzione $f(x)=e^{\sqrt{1-x}}$ è

\begin{enumerate}
\item $1+x+\frac12x^2+\frac16x^3+o(x^3)$
\item $1+x-\frac12x^2+\frac16x^3+o(x^3)$
\item $1+x+\frac12x^2+\frac13x^3+o(x^3)$
\item $e\left(x+\frac12x^2+\frac16x^3+o(x^3)\right)$
\item $e\left(1-\frac12x-\frac1{48}x^3+o(x^3)\right)$
\end{enumerate}
\end{quiz}

% SOLUZIONE: E
% SEZIONE_AM1: 4.1
% SPIEGAZIONE: Sviluppando $\sqrt{1-x}$ e poi l’esponenziale si ottiene $e(1-x/2-x^3/48+o(x^3))$; soluzione calcolata, tabella ufficiale non presente nell’estratto.

% Pag. originale 70/521
\iffalse Pag. originale 70/521 \fi
\begin{quiz}
Sia data la funzione $f(x)=(x+2)^4(3+x)\arctan x$. Sicuramente

\begin{enumerate}
\item $\exists c\in(-3,-2)$ tale che $f(c)=0$
\item $x=-2$ è punto di flesso a tangente orizzontale
\item $x=-3$ è punto critico
\item $\exists c\in(-2,0)$ tale che $f\prime(c)=0$
\item nessuna delle soluzioni precedenti è vera
\end{enumerate}
\end{quiz}

% SOLUZIONE: D
% SEZIONE_AM1: 4.2
% SPIEGAZIONE: Per Rolle, $f(-2)=f(0)=0$ implica un punto $c\in(-2,0)$ con derivata nulla.

% Pag. originale 70/521
\iffalse Pag. originale 70/521 \fi
\begin{quiz}
La funzione $f(x)$, continua in $\mathbb R\setminus\{1,4\}$. Quale delle seguenti situazioni può descrive una funzione con un solo massimo relativo?

\begin{enumerate}
\item $\lim_{x\to\pm\infty}f(x)=+\infty;\ \lim_{x\to1^\pm}f(x)=+\infty;\ \lim_{x\to4^\pm}f(x)=+\infty$
\item $\lim_{x\to\pm\infty}f(x)=-\infty;\ \lim_{x\to1^\pm}f(x)=-\infty;\ \lim_{x\to4^\pm}f(x)=-\infty$
\item $\lim_{x\to\pm\infty}f(x)=+\infty;\ \lim_{x\to1^-}f(x)=-\infty;\ \lim_{x\to4^\pm}f(x)=+\infty$
\item $\lim_{x\to\pm\infty}f(x)=+\infty;\ \lim_{x\to1^\pm}f(x)=+\infty;\ \lim_{x\to4^\pm}f(x)=\pm\infty$
\item $\lim_{x\to\mp\infty}f(x)=\pm\infty;\ \lim_{x\to1^\mp}f(x)=\mp\infty;\ \lim_{x\to4^\pm}f(x)=-\infty$
\end{enumerate}
\end{quiz}

% SOLUZIONE: ?
% SEZIONE_AM1: 4.3
% SPIEGAZIONE: Classificazione basata su estremi relativi; soluzione non riportata nelle pagine fornite. Notazione dei limiti con segni accoppiati da verificare.

% Pag. originale 70/521
\iffalse Pag. originale 70/521 \fi
\begin{quiz}
Sia $z=a+ib$, $a,b\ne0$, $w=3+i$. Allora $u=z\cdot w$ è reale se

\begin{enumerate}
\item $a=-3b$
\item $a=3b$
\item $b=3a$
\item $a=b=0$
\item $a=b=1$
\end{enumerate}
\end{quiz}

% SOLUZIONE: A
% SEZIONE_AM1: 1.5
% SPIEGAZIONE: La parte immaginaria di $(a+ib)(3+i)$ è $a+3b$; si annulla per $a=-3b$.
